\documentclass[12pt,a4paper,oneside]{article}
\usepackage[brazil]{babel}
\usepackage[utf8]{inputenc}
\usepackage{olymp}

\setcounter{problem}{0}

\begin{document}

\begin{problem}{Capivarão lotado}{lotado}{2}

Capivarão é o ônibus que transporta estudantes dentro do Campus da UFV, sendo usado principalmente para deslocamento até o Restaurante Universitário (RU 2), enquanto o RU principal está sendo reformado.

Considere que, em semanas de feriado, quando há menos estudantes, ele aguarda no DCE enquanto há cadeiras vazias.
Em particular, sendo $C$ sua capacidade de pessoas sentadas, ele parte quando o número de estudantes no ônibus é maior ou igual a $C$ (os excedentes vão em pé).
De toda forma, para que os estudantes não aguardem tanto tempo, ele parte após $T$ minutos mesmo se houver cadeira vazia.

Para automatizar o processo, foi instalado um sensor para contabilizar a entrada de estudantes. Sua tarefa é fazer um programa que, dada a quantidade que entra a cada minuto, determina o minuto que ele deve partir,
isto é, quando a capacidade $C$ é atingida ou quanto passam $T$ minutos, o que acontecer primeiro.

%\Notes
%
%Observações, se tiver.

\InputFile

A entrada contém duas linhas: a primeira contém dois números inteiros $C$ e $T$, repectivamente a capacidade do ônibus e o limite de espera;
a segunda linha contém vários números inteiros $Q$, informando sequencialmente quantos estudantes entraram no ônibus a cada minuto.
Restrições: $10 \leq C \leq 50$, $1 \leq T \leq 30$, $0 \leq Q \leq 30$.

\OutputFile

Seu programa deve escrever o minuto que o Capivarão deve partir.

%\Notes

\Example

\begin{example}
\exmp{
20 30
10 5 4 10 1 2
}{
4
}%
\end{example}

\begin{example}
\exmp{
40 5
4 4 4 4 4 4 4 4
}{
5
}%
\end{example}

\begin{example}
\exmp{
30 10
0 1 2 3 4 5 6 7 8 9 10 11 12
}{
9
}%
\end{example}

%Comentários sobre o exemplo, se necessário.
\small
No primeiro exemplo, o ônibus parte após $4$ minutos, que é quando atinge $C=20$ passageiros (entram 10 no minuto 1, mais 5 no minuto 2, mais 4 no minuto 3 e mais 10 no minuto 4).\\
No segundo, ele parte após $5$ minutos, por atingir o tempo limite de espera $T=5$.\\

\end{problem}

\end{document}
